% -*- coding: utf-8 -*-
% 编译方式：XeLaTeX 或 pdfLaTeX（需安装中文字体）
\documentclass[fontset=fandol]{ctexart}

% ==================== 页面与样式 ====================
\usepackage[margin=2.5cm]{geometry}
\usepackage{listings}
\usepackage{xcolor}
\usepackage{hyperref}
\usepackage{amsmath,amssymb}

\hypersetup{
  colorlinks = true,
  linkcolor  = blue,
  urlcolor   = blue,
  pdftitle   = {exactseq 宏包用户手册},
  pdfauthor  = {},
}

% ==================== 代码块样式 ====================
\lstset{
  basicstyle      = \ttfamily\small,
  backgroundcolor = \color{gray!10},
  frame           = single,
  breaklines      = true,
  xleftmargin     = 2em,
  xrightmargin    = 2em,
  gobble          = 2,
  language        = [LaTeX]TeX,
}

% ==================== 命令速记 ====================
\newcommand{\pkg}[1]{\texttt{#1}}
\newcommand{\opt}[1]{\texttt{#1}}
\newcommand{\cmd}[1]{\texttt{\textbackslash #1}}
\newcommand{\marg}[1]{\texttt{\{#1\}}}
\newcommand{\meta}[1]{\textlangle\textit{#1}\textrangle}

% ==================== 文档信息 ====================
\title{\bfseries exactseq 宏包用户手册}
\author{}
\date{v1.0，2026 年 6 月}

\begin{document}
\maketitle
\tableofcontents
\vspace{2em}

\begin{abstract}
\pkg{exactseq} 是一个基于 Ti\textit{k}Z 的 \LaTeX{} 宏包，用于排版数学中的正合列（exact sequence）。
它提供了四条核心命令，支持任意长度的正合列、可配置的箭头样式、最短箭头长度以及标签间距控制。
本宏包使用 \LaTeX3（expl3）语法编写，依赖 \pkg{tikz} 和 \pkg{xparse}。
\end{abstract}

% ================================================================
\section{简介}

在代数学中，正合列是频繁出现的排版需求。传统做法是使用 \cmd{xrightarrow}（来自 \pkg{amsmath}）逐段拼接，
但存在以下问题：

\begin{itemize}
  \item 拼接过程繁琐，容易出错；
  \item 空标签时箭头过短，排版难看（$0 \xrightarrow{} M'$）；
  \item 箭头样式固定为 Computer Modern 字体，无法更换。
\end{itemize}

\pkg{exactseq} 宏包解决了这些问题：

\begin{itemize}
  \item 一条命令即可写出任意长度的正合列；
  \item 基于 Ti\textit{k}Z 绘制箭头，可自由配置箭头样式（Stealth、Latex 等）；
  \item 箭头有最短长度保底，空标签不再过短；
  \item 标签间距可调。
\end{itemize}

% ================================================================
\section{安装与加载}

\subsection{安装}

将 \texttt{exactseq.sty} 放入工作目录，或复制到 \TeX{} 发行版的本地宏包路径：

\begin{lstlisting}
  # macOS / Linux
  cp exactseq.sty ~/texmf/tex/latex/exactseq/
  texhash ~/texmf

  # Windows
  copy exactseq.sty %USERPROFILE%\texmf\tex\latex\exactseq\
  texhash %USERPROFILE%\texmf
\end{lstlisting}

\subsection{加载}

\begin{lstlisting}
  \usepackage{exactseq}
\end{lstlisting}

本宏包自动加载 \pkg{tikz}（含 \texttt{arrows.meta} 与 \texttt{cd} 库）和 \pkg{xparse}。

% ================================================================
\section{用户命令}

\subsection{\cmd{exactseq} —— 任意长度正合列}

这是最核心的命令，通过递归解析括号对实现任意长度序列的拼接。

\medskip
\textbf{语法}
\begin{lstlisting}
  \exactseq{<首个对象>}[<映射1>]{<对象2>}[<映射2>]{<对象3>}...
\end{lstlisting}

\medskip
\textbf{规则}
\begin{itemize}
  \item 第一个参数 \marg{对象} 无前置箭头，直接输出；
  \item 之后每个 \texttt{[}\meta{映射}\texttt{]\{}\meta{对象}\texttt{\}} 对输出
        ``$\xrightarrow{\text{映射}}$ 对象''；
  \item \texttt{[}\texttt{]} 为空时输出无标签箭头（使用最短长度）；
  \item 序列以 \texttt{\{}\meta{最后一个对象}\texttt{\}} 结尾。
\end{itemize}

\medskip
\textbf{示例}
\begin{lstlisting}
  % 短正合列
  \exactseq{0}[]{M'}[f]{M}[g]{M''}[]{0}
  % 输出：0 -> M' -f-> M -g-> M'' -> 0

  % 多段序列
  \exactseq{\cdots}[]{M_{i-1}}[]{M_i}[]{M_{i+1}}[]{\cdots}

  % 长标签自动伸展
  \exactseq{0}[]{M \otimes_R A}[1_M \otimes f]{M \otimes_R B}[]{0}
\end{lstlisting}

\subsection{\cmd{ses} —— 短正合列}

预置好 ``$0 \to \cdot \to \cdot \to \cdot \to 0$'' 的格式，适合短正合列。

\medskip
\textbf{语法}
\begin{lstlisting}
  \ses[<映射1>][<映射2>]{<M'>}{<M>}{<M''>}
\end{lstlisting}

\medskip
\textbf{示例}
\begin{lstlisting}
  \ses[f][g]{M'}{M}{M''}
  % 输出：0 -> M' -f-> M -g-> M'' -> 0
\end{lstlisting}

\subsection{\cmd{injseq} —— 左侧正合列}

预置好 ``$0 \to \cdot \to \cdot$'' 的格式，适合表示单射（injection）。

\medskip
\textbf{语法}
\begin{lstlisting}
  \injseq[<映射>]{<M'>}{<M>}
\end{lstlisting}

\medskip
\textbf{示例}
\begin{lstlisting}
  \injseq[f]{M'}{M}
  % 输出：0 -> M' -f-> M
\end{lstlisting}

\subsection{\cmd{surjseq} —— 右侧正合列}

预置好 ``$\cdot \to \cdot \to 0$'' 的格式，适合表示满射（surjection）。

\medskip
\textbf{语法}
\begin{lstlisting}
  \surjseq[<映射>]{<M>}{<M''>}
\end{lstlisting}

\medskip
\textbf{示例}
\begin{lstlisting}
  \surjseq[g]{M}{M''}
  % 输出：M -g-> M'' -> 0
\end{lstlisting}

% ================================================================
\section{配置选项}

使用 \cmd{exactseqset} 命令在导言区或正文中调整各项参数。

\medskip
\textbf{语法}
\begin{lstlisting}
  \exactseqset{ <键1> = <值1>, <键2> = <值2>, ... }
\end{lstlisting}

\medskip
\textbf{所有可配置选项}

\begin{center}
\begin{tabular}{lllp{6cm}}
  \hline
  键 & 类型 & 默认值 & 说明 \\
  \hline
  \opt{arrow style} & 字串 & \texttt{->} &
    箭头样式，支持所有 Ti\textit{k}Z \texttt{arrows.meta} 样式 \\
  \opt{min length}  & 尺寸 & \texttt{1.8em} &
    箭头最短长度（标签长于此值时自动伸展） \\
  \opt{label style} & 字串 & \texttt{\textbackslash scriptstyle} &
    标签的字体命令 \\
  \opt{label sep}   & 尺寸 & \texttt{1.5pt} &
    标签与箭头线的垂直间距 \\
  \hline
\end{tabular}
\end{center}

\medskip
\textbf{示例}
\begin{lstlisting}
  % 更换箭头为 Stealth 样式
  \exactseqset{arrow style={-{Stealth}}}

  % 自定义尺寸的 Latex 箭头 + 更宽的最短长度
  \exactseqset{
    arrow style = {-{Latex[length=3.5pt, width=2pt]}},
    min length  = 2.2em,
  }

  % 调整标签间距
  \exactseqset{label sep = 2.5pt}
\end{lstlisting}

% ================================================================
\section{箭头样式参考}

本宏包使用 Ti\textit{k}Z 的 \texttt{arrows.meta} 库，以下列出常用的数学场景箭头样式。

\begin{center}
\begin{tabular}{ll}
  \hline
  \opt{arrow style} 值 & 效果描述 \\
  \hline
  \texttt{->} & 默认实心三角，接近 CM 风格 \\
  \texttt{-\{Stealth\}} & 实心三角，线条更粗 \\
  \texttt{-\{Latex\}} & 类似 CM 但更均匀 \\
  \texttt{-\{Latex[length=3.5pt,width=2pt]\}} & Latex 箭头，自定义尺寸 \\
  \texttt{-\{Computer Modern Rightarrow\}} & 最接近 CM 字体的原始箭头 \\
  \texttt{->>} & 双线箭头，适合表示满射 \\
  \texttt{-\{>\}|>\}} & 双箭头（两个并列箭头） \\
  \texttt{-\{Hooks\}} & 钩子，适合表示单射 \\
  \texttt{-\{Straight Barb\}} & 细长直箭头 \\
  \texttt{-\{Bar\}} & 竖线结尾（无箭头） \\
  \hline
\end{tabular}
\end{center}

更多样式请参阅 Ti\textit{k}Z 手册第 16 章（命令行运行 \texttt{texdoc tikz}）。

% ================================================================
\section{完整示例}

\subsection{基础用法}

\begin{lstlisting}
  \documentclass{article}
  \usepackage{amsmath}
  \usepackage{exactseq}

  \begin{document}

  行内正合列：
  $\exactseq{0}[]{M'}[f]{M}[g]{M''}[]{0}$

  行间正合列：
  \[
    \exactseq{0}[]{M_1}[f_1]{M_2}[f_2]{M_3}[]{\cdots}[]{\cdots}
  \]

  短正合列：
  \[
    \ses[f][g]{M'}{M}{M''}
  \]

  \end{document}
\end{lstlisting}

\subsection{自定义箭头样式}

\begin{lstlisting}
  \documentclass{article}
  \usepackage{exactseq}

  % 全局设置
  \exactseqset{
    arrow style = {-{Stealth}},
    min length  = 2em,
  }

  \begin{document}

  \[
    \exactseq{0}[]{\ker f}[i]{M}[f]{N}[\pi]{\coker f}[]{0}
  \]

  % 局部切换为双线箭头（满射用）
  \exactseqset{arrow style={->>}}
  \[
    \exactseq{M}[g]{M''}[]{0}
  \]

  \end{document}
\end{lstlisting}

% ================================================================
\section{注意事项}

\begin{enumerate}
  \item 所有命令仅可在数学模式（\texttt{\$...\$} 或 \texttt{\textbackslash[...\textbackslash]}）中使用。
  \item 箭头标签会以 \cmd{scriptstyle} 排版（与 \cmd{xrightarrow} 行为一致）。
  \item \cmd{exactseq} 的第一个参数是无前置箭头的对象，因此以 \texttt{0} 或 \texttt{\textbackslash cdots}
        开始时，其后必须跟 \texttt{[]\{\}} 对。
  \item 若序列过长导致行间溢出，可手动换行（本宏包不处理自动换行）。
\end{enumerate}

% ================================================================
\section{版本历史}

\begin{description}
  \item[v1.0]（2026-06-01）
    \begin{itemize}
      \item 初始版本。
      \item 提供 \cmd{exactseq}、\cmd{ses}、\cmd{injseq}、\cmd{surjseq} 四条命令。
      \item 基于 Ti\textit{k}Z 的箭头绘制，支持 \texttt{arrows.meta} 样式。
      \item 可配置最短箭头长度、标签间距。
    \end{itemize}
\end{description}

\end{document}
